\clearpage
\fancyhead[L]{\small DIAMETER OF THE SUN.}
\hypertarget{native-530-lower}{}\label{native-530-lower}
\begin{center}\small Göttingische gelehrte Anzeigen. Stück 95. pp. 937..940. June 16, 1810.\\[3pt]\rule{.72\linewidth}{.3pt}\end{center}

In no. 113 of last year's \emph{Göttingische gelehrte Anzeigen} [p. 517 of this volume], we reported the remarkable results that Mr. \textsc{von Lindenau} had obtained for the Sun's diameter from the Greenwich observations of 1750... 1786. The astronomer has since extended his investigations to the more recent observations of 1787... 1798 and has communicated the results to our Professor \textsc{Gauss} in a letter of 23 May. It is interesting that they agree very well with the earlier results, and they therefore deserve all the more attention, at least from a practical point of view. First of all, the semidiameters reduced to the mean distance again show a kind of periodic inequality at different seasons, as the following table illustrates.

\begin{center} \begin{tabular}{l|c||l|c} & Solar semidiameter & & Solar semidiameter \\ \hline January & $959''70$ & July & $959''22$ \\ February & $959.99$ & August & $959.98$ \\ March & $960.41$ & September & $960.19$ \\ April & $959.80$ & October & $960.10$ \\ May & $959.81$ & November & $960.07$ \\ June & $959.00$ & December & $958.75$ \\ \end{tabular} \end{center}

These observations can be represented quite well by setting $\frac{bb}{aa} - 1 = \frac{1}{92}$ in the formula on the page cited, p. 1125 [p. 519 of this volume] (where, incidentally, a misprint gives 0.9631 instead of 0.09631). This value is somewhat smaller than the one found there, but is also positive.

Mr. \textsc{von Lindenau} finds the mean solar semidiameters from all the individual observations in the different years as follows:

\begin{center} \begin{tabular}{l|c|c} & \begin{tabular}{c}Solar\\semidiameter\end{tabular} & \begin{tabular}{c}Number of\\observations\end{tabular} \\ \hline 1787 & $959''26$ & 62 \\ 1788 & $959.10$ & 75 \\ 1789 & $959.68$ & 95 \\ 1790 & $959.33$ & 80 \\ 1791 & $959.51$ & 91 \\ 1792 & $960.24$ & 65 \\ 1793 & $959.59$ & 75 \\ 1794 & $959.96$ & 62 \\ 1795 & $960.05$ & 64 \\ 1796 & $960.02$ & 48 \\ 1797 & $959.99$ & 70 \\ 1798 & $960.03$ & 72 \\ \end{tabular} \end{center}

\SourceLeaf{531}
As to the apparent progressive decrease in the solar semidiameter, grouping \textsc{Maskelyne}'s observations over thirty-three years into three periods gives the following results:

\begin{center} \begin{tabular}{c|c} & Solar semidiameter \\ \hline 1765... 1775 & $961''66$ \\ 1776... 1786 & $960.22$ \\ 1787... 1798 & $959.77$ \\ \end{tabular} \end{center}

From 116 horizontal solar semidiameters observed by \textsc{Piazzi} in 1791... 1793 Mr. \textsc{von Lindenau} obtains $961''21$, and from \textsc{Bradley}'s observations of 1755... 1760 he obtains $961''86$. The close agreement of measurements from such different times confirms that the size of the Sun remains constant. Mr. \textsc{von Lindenau}, however, attributes the progressive decrease over the long period of \textsc{Maskelyne}'s observations to a change in the observer's eyesight, with irradiation apparently becoming less as he grew older. Mr. \textsc{von Lindenau} himself calls this only a conjecture that is not unworthy of further investigation.

The considerable difference between the horizontal and vertical semidiameters appears again in the more recent observations. Mr. \textsc{von Lindenau}'s calculations on this point give the following results:

\begin{center}\small\setlength{\tabcolsep}{3pt} \begin{tabular}{l|c|c||l|c|c} & \begin{tabular}{c}Vertical\\semidiameter\end{tabular} & \begin{tabular}{c}Number of\\observations\end{tabular} & & \begin{tabular}{c}Vertical\\semidiameter\end{tabular} & \begin{tabular}{c}Number of\\observations\end{tabular} \\ \hline 1787 & $962''42$ & 72 & 1793 & $963''22$ & 49 \\ 1788 & $962.44$ & 102 & 1794 & $963.13$ & 72 \\ 1789 & $963.10$ & 91 & 1795 & $963.22$ & 60 \\ 1790 & $963.03$ & 66 & 1796 & $962.40$ & 61 \\ 1791 & $963.10$ & 57 & 1797 & $962.25$ & 50 \\ 1792 & $963.30$ & 48 & 1798 & $962.10$ & 53 \\ \end{tabular} \end{center}

The arithmetic mean of them all is $962''82$, which is $2''96$ greater than the mean of the horizontal semidiameters. The observations of 1765... 1786 gave $2''80$ for this difference. From 71 observations by \textsc{Piazzi}, Mr. \textsc{von Lindenau} finds that the mean vertical semidiameter is $963''01$ and that the horizontal one is $1''80$ smaller. Observations by \textsc{Bouvard} also gave Mr. \textsc{von Lindenau} a similar difference.

This agreement among the results from three different observers is indeed very striking. If, however, other considerations carry enough weight to make it difficult to accept a prolate ellipsoidal shape for the Sun, one might be more inclined to look elsewhere for the reason for this difference. Perhaps the most appropriate place to look is the way the observations are made. In measuring zenith distances one usually brings the limb of the Sun into contact with the horizontal wire \emph{from the outside}. The resulting zenith distance of the upper limb must therefore be too small by half the thickness of the wire, and that of the lower limb too large by the same amount. Solar diameters determined in this way must accordingly be too large by the full thickness of the wire.
